\documentclass[11pt,letterpaper]{article}

% =============================================================================
% CONFIGURACIÓN DE PAQUETES Y TIPOGRAFÍA
% =============================================================================
\usepackage[utf8]{inputenc}
\usepackage[T1]{fontenc}
\usepackage[spanish,es-tabla]{babel}
\usepackage{lmodern}
\usepackage[margin=2.3cm,top=2.5cm,bottom=2.5cm]{geometry}
\usepackage{amsmath,amssymb}
\usepackage{graphicx}
\usepackage{booktabs,tabularx,longtable,array,multirow}
\usepackage{xcolor}
\usepackage{listings}
\usepackage{tikz}
\usetikzlibrary{shapes.geometric,arrows.meta,positioning,calc,fit,backgrounds}
\usepackage{fancyhdr}
\usepackage{enumitem}
\usepackage{titlesec}
\usepackage{microtype}
\usepackage[hidelinks,pdftex,bookmarks=true,bookmarksnumbered=true]{hyperref}

% =============================================================================
% PALETA DE COLORES INSTITUCIONAL
% =============================================================================
\definecolor{pujblue}{RGB}{0, 43, 91}       % Azul Javeriana
\definecolor{pujgold}{RGB}{181, 133, 0}     % Dorado Javeriana
\definecolor{accentdark}{RGB}{20, 30, 45}
\definecolor{codebg}{RGB}{248, 249, 250}
\definecolor{codeframe}{RGB}{220, 224, 230}
\definecolor{sqlkeyword}{RGB}{0, 70, 160}
\definecolor{sqlstring}{RGB}{160, 40, 40}
\definecolor{sqlcomment}{RGB}{100, 110, 120}

% =============================================================================
% ESTILO DE CÓDIGO (LISTINGS)
% =============================================================================
\lstdefinestyle{sqlstyle}{
    language=SQL,
    backgroundcolor=\color{codebg},
    basicstyle=\ttfamily\footnotesize,
    keywordstyle=\color{sqlkeyword}\bfseries,
    stringstyle=\color{sqlstring},
    commentstyle=\color{sqlcomment}\itshape,
    breaklines=true,
    breakatwhitespace=true,
    showstringspaces=false,
    frame=single,
    rulecolor=\color{codeframe},
    numbers=left,
    numberstyle=\tiny\color{gray},
    numbersep=7pt,
    tabsize=2,
    morekeywords={
        SERIAL, BYTEA, INTERVAL, TIMESTAMPTZ, tsrange, GREATEST, date_trunc,
        FILTER, DO, PERFORM, STRICT, FOR, EACH, ROW, EXECUTE, FUNCTION,
        PARTITION, OF, ATTACH, TABLESAPCE, TABLESPACE
    }
}
\lstset{style=sqlstyle}

% =============================================================================
% ENCABEZADOS Y PIES DE PÁGINA
% =============================================================================
\pagestyle{fancy}
\fancyhf{}
\renewcommand{\headrulewidth}{0.6pt}
\renewcommand{\footrulewidth}{0.4pt}
\fancyhead[L]{\small\scshape\color{pujblue}Pontificia Universidad Javeriana}
\fancyhead[R]{\small\color{gray}Proyecto Restaurante | BD 2026-3}
\fancyfoot[L]{\footnotesize\color{gray}Ingeniería de Sistemas - Bases de Datos}
\fancyfoot[R]{\thepage}

% Títulos de secciones estilizados
\titleformat{\section}{\color{pujblue}\normalfont\Large\bfseries}{\thesection.}{0.8em}{}[\titlerule]
\titleformat{\subsection}{\color{accentdark}\normalfont\large\bfseries}{\thesubsection.}{0.6em}{}
\titleformat{\subsubsection}{\color{pujgold!80!black}\normalfont\normalsize\bfseries}{\thesubsubsection.}{0.5em}{}

% =============================================================================
% DOCUMENTO PRINCIPAL
% =============================================================================
\begin{document}

% -----------------------------------------------------------------------------
% PORTADA INSTITUCIONAL
% -----------------------------------------------------------------------------
\begin{titlepage}
    \centering
    \vspace*{0.5cm}
    
    {\scshape\LARGE Pontificia Universidad Javeriana \par}
    \vspace{0.4cm}
    {\scshape\large Facultad de Ingeniería \par}
    {\scshape\large Departamento de Ingeniería de Sistemas \par}
    \vspace{0.3cm}
    {\large Bases de Datos: Administración y Aplicaciones (Período 2026-3) \par}
    
    \vspace{1.8cm}
    \rule{\linewidth}{1.5pt} \\[0.4cm]
    {\huge\bfseries\color{pujblue} PROYECTO RESTAURANTE \\[0.25cm] \large Diseño Relacional, Almacenamiento Físico, Integridad Avanzada y Consultas Analíticas en PostgreSQL 18 \par}
    \rule{\linewidth}{1.5pt} \\[1.8cm]
    
    \begin{minipage}[t]{0.48\textwidth}
        \textbf{Presentado por:} \\[0.2cm]
        Estudiante(s): \\
        \textit{García \& Pico} \\
        Ingeniería de Sistemas
    \end{minipage}
    \hfill
    \begin{minipage}[t]{0.48\textwidth}
        \raggedleft
        \textbf{Docente Titular:} \\[0.2cm]
        \textbf{Ing. Ricardo Naranjo Faccini, MSc.} \\
        Profesor Asistente \\
        Departamento de Ingeniería de Sistemas
    \end{minipage}
    
    \vfill
    {\large Bogotá D.C., Colombia \\ Noviembre de 2026 \par}
\end{titlepage}

% -----------------------------------------------------------------------------
% SUMARIO Y TABLA DE CONTENIDO
% -----------------------------------------------------------------------------
\pagenumbering{roman}
\setcounter{page}{1}

\section*{Resumen Ejecutivo}
El presente documento contiene el informe técnico y formal de entrega correspondiente al \textbf{Proyecto Restaurante} para la asignatura de \textit{Bases de Datos: Administración y Aplicaciones} de la Pontificia Universidad Javeriana. 

El objetivo primordial consiste en diseñar, implementar y documentar una base de datos relacional robusta en \textbf{PostgreSQL 18} para soportar la operación integral de un restaurante de alta demanda. El sistema gestiona mesas, reservaciones, menús, pedidos y entregas, garantizando reglas de negocio críticas: control estricto del cupo total de mesas, prevención matemática de cruces de reservaciones en ventanas mínimas de dos horas, despacho de platos hacia cocineros autorizados según especialidades capacitadas, y trazabilidad completa del ciclo de vida de los pedidos para los meseros y maitres.

Asimismo, se expone la arquitectura física basada en \textbf{tablespaces} en directorios independientes de almacenamiento, particionamiento de tablas por rangos y fechas, disparadores (\textit{triggers}) de integridad referencial activa, y el catálogo exhaustivo de consultas SQL optimizadas para resolver los casos de uso operativos y reportes analíticos de gestión.

\vspace{0.8cm}
\tableofcontents
\newpage
\listoftables
\newpage

% -----------------------------------------------------------------------------
% CUERPO DEL INFORME TÉCNICO
% -----------------------------------------------------------------------------
\pagenumbering{arabic}
\setcounter{page}{1}

\section{Descripción General y Roles de Usuario}

\subsection{Contexto y Objetivos del Sistema}
El proyecto aborda la problemática operacional de establecimientos gastronómicos en donde confluyen simultáneamente múltiples perfiles de usuarios y flujos transaccionales altamente concurrentes. El modelo debe garantizar:
\begin{itemize}[leftmargin=*]
    \item \textbf{Aforo y Asignación de Mesas:} No exceder en ningún momento el cupo máximo de comensales del local ni solapar reservas sobre un mismo recurso físico.
    \item \textbf{Gestión de Horarios y Reservaciones:} Todo turno de reserva estima una ocupación predeterminada de dos (2) horas, calculando la disponibilidad a través de álgebra de intervalos (\texttt{tsrange}).
    \item \textbf{Especialización en Cocina:} Cada cocinero solo puede visualizar, tomar y marcar como preparados aquellos platos para los cuales ha sido formalmente acreditado mediante una especialidad.
    \item \textbf{Gestión Operativa de Pedidos:} Los meseros registran órdenes asociadas a mesas y clientes, gestionando la transición de estados (\textit{Solicitada $\rightarrow$ En preparación $\rightarrow$ Preparada $\rightarrow$ Entregada}).
    \item \textbf{Inteligencia de Negocio:} Generación dinámica de reportes consolidados por día, semana y mes para ventas, afluencia de reservas y platos de mayor demanda.
\end{itemize}

\subsection{Matriz de Roles y Responsabilidades}
\begin{table}[h!]
\centering
\small
\begin{tabularx}{\textwidth}{l p{4.2cm} X}
\toprule
\textbf{Rol} & \textbf{Alcance Operativo} & \textbf{Acciones Principales en Base de Datos} \\
\midrule
\textbf{Administrador} & Control maestro de infraestructura, personal y finanzas & Mantenimiento de mesas (CRUD), actualización de carta/precios, asignación de roles y especialidades a empleados, y consulta de reportes analíticos. \\
\addlinespace
\textbf{Maitre} & Control de sala, aforo y relaciones de acceso & Registro, consulta y cancelación de reservaciones; verificación de cupo total en tiempo real; atención de alertas tempranas ($\le 2$ horas). \\
\addlinespace
\textbf{Mesero} & Toma de comandas y servicio a mesa & Creación de pedidos con múltiples ítems, consulta de comandas listas en cocina y registro de entrega efectiva en mesa. \\
\addlinespace
\textbf{Cocinero} & Elaboración culinaria en cocina & Visualización de cola de comandas autorizadas priorizadas por tiempo descendente de cocción, toma de órdenes y notificación de plato listo. \\
\addlinespace
\textbf{Cliente} & Comensal (interacción indirecta) & Titular de reservaciones y pedidos; consulta de historial de consumos y visitas. \\
\bottomrule
\end{tabularx}
\caption{Matriz de roles y permisos del sistema del restaurante.}
\label{tab:roles}
\end{table}

\newpage
\section{Casos de Uso del Sistema}

A continuación se detallan los siete casos de uso troncales y las funcionalidades transversales estipuladas en la especificación oficial:

\begin{enumerate}[label=\textbf{CU-\arabic*:},leftmargin=*]
    \item \textbf{Gestión de Mesas (Administrador):} Creación, actualización y retiro de mesas físicas, asignando su número único y número de asientos disponibles. Cálculo automático del aforo o cupo total mediante $\sum \text{sillas}$.
    \item \textbf{Gestión de Reservaciones (Maitre):} Asignación de mesa y horario a clientes. Se exige validar que la mesa seleccionada soporte el número de comensales, que no existan traslapes temporales con otras reservas en dicha mesa (considerando duración estándar de 2 horas), y que la suma de asistentes concurrentes no sobrepase el cupo del restaurante.
    \item \textbf{Gestión del Menú (Administrador):} Catálogo de ítems clasificados en cuatro categorías obligatorias (\textit{Entrada, Plato Fuerte, Postre, Bebida}). Cada plato registra nombre, descripción textual, precio unitario vigente y tiempo estimado de preparación (\texttt{INTERVAL}). Registro de especialidades por cocinero.
    \item \textbf{Registro de Pedidos (Mesero):} Creación de pedidos cabecera vinculados a una mesa y un cliente. Inclusión de una o múltiples líneas de orden (\texttt{ordenes}) capturando el precio unitario vigente en ese instante. Estado inicial automático en \textit{Solicitado/En preparación}.
    \item \textbf{Preparación de Platos (Cocinero):} Vista en cocina donde cada cocinero accede únicamente a las órdenes de platos que está capacitado para elaborar. Las órdenes se ordenan de forma descendente por tiempo de cocción (\textit{los platos más demorados de primero}) y secundariamente por antigüedad de solicitud. Al culminar, actualiza el estado a \textit{Preparado}.
    \item \textbf{Gestión de Entregas (Mesero):} Monitoreo de ítems con estado \textit{Preparado} listos para servir. El mesero efectúa la entrega material a la mesa y actualiza el estado del ítem a \textit{Entregado}.
    \item \textbf{Reportes y Estadísticas (Administrador):} Consultas analíticas agrupadas dinámicamente por día, semana o mes (\texttt{date\_trunc}) que consolidan:
    \begin{itemize}
        \item Volumen y estados de reservaciones.
        \item Ranking de los platos más solicitados (unidades e ingresos).
        \item Facturación global y ventas verificadas por fecha de entrega.
    \end{itemize}
\end{enumerate}

\paragraph{Funcionalidades Adicionales Implementadas:}
\begin{itemize}
    \item \textbf{Sistema de Alertas y Notificaciones:} Alertas automáticas para el Maitre ante reservas que inician en $\le 2$ horas, y para el Mesero ante órdenes listas para entrega.
    \item \textbf{Historial Completo de Clientes:} Métrica consolidada por cliente con visitas históricas, pedidos consumidos y valor monetario acumulado.
    \item \textbf{Gestión Integral de Personal:} Control de altas, bajas lógicas (\texttt{activo}) y modificación de roles múltiples mediante la entidad asociativa \texttt{actuaciones}.
\end{itemize}

\newpage
\section{Diseño de la Base de Datos y Modelo Entidad-Relación}

\subsection{Modelo Entidad-Relación Obligatorio}
El sistema implementa de manera estricta el esquema relacional de 12 entidades especificado por la cátedra. La Figura~\ref{fig:diagrama_er} ilustra la topología relacional y las dependencias foráneas.

\begin{figure}[h!]
\centering
\begin{tikzpicture}[
    scale=0.85, every node/.style={transform shape},
    entity/.style={rectangle, draw=pujblue, fill=blue!5, thick, rounded corners=2pt, minimum width=3.2cm, inner sep=6pt},
    rel/.style={-{Latex[length=2.5mm]}, thick, pujblue!80!black},
    textlabel/.style={font=\scriptsize\ttfamily}
]
    % Filas del diagrama
    % Nivel 1: Menú y Tipos
    \node[entity] (tipos) at (0, 8) {
        \textbf{Tipos} \\ \hline
        \scriptsize \underline{id}: SERIAL \\
        \scriptsize nombre: TEXT
    };
    
    \node[entity] (platos) at (5, 8) {
        \textbf{Platos} \\ \hline
        \scriptsize \underline{id}: SERIAL \\
        \scriptsize tipo\_id: INT (FK) \\
        \scriptsize nombre: TEXT \\
        \scriptsize tiempo: INTERVAL \\
        \scriptsize precio: NUMERIC
    };
    
    \node[entity] (ordenes) at (11, 8) {
        \textbf{Ordenes} \\ \hline
        \scriptsize \underline{id}: SERIAL \\
        \scriptsize plato\_id: INT (FK) \\
        \scriptsize pedido\_id: INT (FK) \\
        \scriptsize estado: INT \\
        \scriptsize cantidad: INT \\
        \scriptsize solicitado: TIMESTAMPTZ
    };

    % Nivel 2: Especialidades y Preparaciones
    \node[entity] (esp) at (2.5, 4.5) {
        \textbf{Especialidades} \\ \hline
        \scriptsize \underline{id}: SERIAL \\
        \scriptsize cocinero\_id: INT (FK) \\
        \scriptsize plato\_id: INT (FK)
    };
    
    \node[entity] (prep) at (7.5, 4.5) {
        \textbf{Preparaciones} \\ \hline
        \scriptsize \underline{id}: SERIAL \\
        \scriptsize cocinero\_id: INT (FK) \\
        \scriptsize orden\_id: INT (FK)
    };
    
    \node[entity] (pedidos) at (11, 4.5) {
        \textbf{Pedidos} \\ \hline
        \scriptsize \underline{id}: SERIAL \\
        \scriptsize cliente\_id: INT (FK) \\
        \scriptsize mesero\_id: INT (FK)
    };

    % Nivel 3: Usuarios, Roles, Actuaciones
    \node[entity] (act) at (11, 1) {
        \textbf{Actuaciones} \\ \hline
        \scriptsize \underline{id}: SERIAL \\
        \scriptsize rol\_id: INT (FK) \\
        \scriptsize usuario\_id: INT (FK)
    };

    \node[entity] (roles) at (11, -2) {
        \textbf{Roles} \\ \hline
        \scriptsize \underline{id}: SERIAL \\
        \scriptsize nombre: TEXT
    };

    \node[entity] (usuarios) at (5.5, 1) {
        \textbf{Usuarios} \\ \hline
        \scriptsize \underline{id}: SERIAL \\
        \scriptsize nombre: TEXT \\
        \scriptsize clave: BYTEA \\
        \scriptsize fecha\_clave: TIMESTAMP
    };

    % Nivel 4: Reservaciones, Horarios, Mesas
    \node[entity] (res) at (0.5, 1) {
        \textbf{Reservaciones} \\ \hline
        \scriptsize \underline{id}: SERIAL \\
        \scriptsize cliente\_id: INT (FK) \\
        \scriptsize cantidad: INT \\
        \scriptsize estado: INT
    };

    \node[entity] (horarios) at (3.5, -2) {
        \textbf{Horarios} \\ \hline
        \scriptsize \underline{id}: SERIAL \\
        \scriptsize mesa\_id: INT (FK) \\
        \scriptsize reservacion\_id: INT (FK) \\
        \scriptsize inicio: TIMESTAMP \\
        \scriptsize duracion: INTERVAL
    };

    \node[entity] (mesas) at (8, -2) {
        \textbf{Mesas} \\ \hline
        \scriptsize \underline{id}: SERIAL \\
        \scriptsize sillas: INT
    };

    % Conexiones relacionales
    \draw[rel] (tipos) -- (platos) node[midway,above,textlabel]{1:N};
    \draw[rel] (platos) -- (ordenes) node[midway,above,textlabel]{1:N};
    \draw[rel] (platos) -- (esp) node[midway,right,textlabel]{1:N};
    \draw[rel] (ordenes) -- (prep) node[midway,above,textlabel]{1:N};
    \draw[rel] (pedidos) -- (ordenes) node[midway,right,textlabel]{1:N};
    
    \draw[rel] (usuarios) -- (esp) node[midway,left,textlabel]{1:N};
    \draw[rel] (usuarios) -- (prep) node[midway,right,textlabel]{1:N};
    \draw[rel] (usuarios) -- (pedidos) node[midway,above,textlabel]{1:N};
    \draw[rel] (usuarios) -- (res) node[midway,above,textlabel]{1:N};
    \draw[rel] (usuarios) -- (act) node[midway,above,textlabel]{1:N};
    
    \draw[rel] (roles) -- (act) node[midway,right,textlabel]{1:N};
    \draw[rel] (res) -- (horarios) node[midway,left,textlabel]{1:N};
    \draw[rel] (mesas) -- (horarios) node[midway,above,textlabel]{1:N};

\end{tikzpicture}
\caption{Diagrama Entidad-Relación formal obligatorio del sistema del restaurante.}
\label{fig:diagrama_er}
\end{figure}

\subsection{Normalización del Esquema}
El diseño satisface la \textbf{Tercera Forma Normal (3FN)} y la \textbf{Forma Normal de Boyce-Codd (FNBC)}:
\begin{enumerate}
    \item \textbf{1FN:} Todos los atributos son atómicos y no existen grupos repetitivos. Los roles múltiples de los usuarios se resuelven a través de la tabla asociativa \texttt{actuaciones}.
    \item \textbf{2FN:} Toda columna que no forma parte de una clave candidata depende por completo de la clave primaria completa (claves subrogadas simples mediante secuencias \texttt{SERIAL}).
    \item \textbf{3FN:} No existen dependencias transitivas entre atributos no clave. Por ejemplo, el precio y categoría no se embeben en la orden, sino que se normalizan en \texttt{platos} y \texttt{tipos}, guardando en la orden únicamente el snapshot transaccional (\texttt{precio\_unitario}).
\end{enumerate}

\newpage
\section{Diccionario de Datos Exhaustivo}

A continuación se documenta la estructura de campos, tipos nativos de PostgreSQL 18, claves y restricciones de integridad para las 12 entidades del sistema:

% 1. Usuarios, Roles y Actuaciones
\subsection{Módulo de Usuarios, Seguridad y Roles}

\begin{table}[h!]
\centering
\scriptsize
\begin{tabularx}{\textwidth}{l l l p{2.5cm} X}
\toprule
\textbf{Columna} & \textbf{Tipo Dato} & \textbf{Nulo} & \textbf{Claves / Constraints} & \textbf{Descripción Funcional} \\
\midrule
\multicolumn{5}{l}{\textbf{Tabla: public.usuarios} (Usuarios del sistema: clientes, cocineros, meseros, maitres, admin)} \\
\texttt{id} & \texttt{SERIAL} & No & \texttt{PRIMARY KEY} & Identificador único subrogado del usuario. \\
\texttt{nombre} & \texttt{TEXT} & No & -- & Nombre completo del usuario o razón social. \\
\texttt{clave} & \texttt{BYTEA} & No & -- & Hash criptográfico de autenticación (SHA-256). \\
\texttt{fecha\_clave} & \texttt{TIMESTAMP} & No & \texttt{DEFAULT now()} & Marca temporal del último cambio de contraseña. \\
\texttt{activo} & \texttt{BOOLEAN} & No & \texttt{DEFAULT true} & Estado de habilitación operativa para login. \\
\midrule
\multicolumn{5}{l}{\textbf{Tabla: public.roles} (Catálogo de roles del restaurante)} \\
\texttt{id} & \texttt{SERIAL} & No & \texttt{PRIMARY KEY} & Identificador único del rol. \\
\texttt{nombre} & \texttt{TEXT} & No & \texttt{UNIQUE} & Nombre del rol: Administrador, Maitre, Mesero, Cocinero, Cliente. \\
\midrule
\multicolumn{5}{l}{\textbf{Tabla: public.actuaciones} (Asociación N:M entre usuarios y roles)} \\
\texttt{id} & \texttt{SERIAL} & No & \texttt{PRIMARY KEY} & Identificador de asignación de rol. \\
\texttt{rol\_id} & \texttt{INTEGER} & No & \texttt{FK -> roles(id)} & Rol concedido al usuario. \\
\texttt{usuario\_id} & \texttt{INTEGER} & No & \texttt{FK -> usuarios(id)} & Usuario al cual se le otorga el rol. \\
\bottomrule
\end{tabularx}
\caption{Diccionario de datos: Usuarios, Roles y Actuaciones.}
\label{tab:dict_usuarios}
\end{table}

\subsection{Módulo de Menú y Especialidades}

\begin{table}[h!]
\centering
\scriptsize
\begin{tabularx}{\textwidth}{l l l p{2.8cm} X}
\toprule
\textbf{Columna} & \textbf{Tipo Dato} & \textbf{Nulo} & \textbf{Claves / Constraints} & \textbf{Descripción Funcional} \\
\midrule
\multicolumn{5}{l}{\textbf{Tabla: public.tipos} (Categorías del menú)} \\
\texttt{id} & \texttt{SERIAL} & No & \texttt{PRIMARY KEY} & Identificador de categoría. \\
\texttt{nombre} & \texttt{TEXT} & No & \texttt{UNIQUE} & Entrada, Plato Fuerte, Postre o Bebida. \\
\midrule
\multicolumn{5}{l}{\textbf{Tabla: public.platos} (Catálogo de platos y bebidas ofrecidos)} \\
\texttt{id} & \texttt{SERIAL} & No & \texttt{PRIMARY KEY} & Identificador del plato. \\
\texttt{tipo\_id} & \texttt{INTEGER} & No & \texttt{FK -> tipos(id)} & Categoría a la que pertenece el ítem. \\
\texttt{nombre} & \texttt{TEXT} & No & -- & Nombre comercial del plato o bebida. \\
\texttt{descripcion} & \texttt{TEXT} & Sí & -- & Reseña de ingredientes y presentación. \\
\texttt{tiempo} & \texttt{INTERVAL} & No & \texttt{CHECK (> '0')} & Duración promedio estimada de cocción/preparación. \\
\texttt{precio} & \texttt{NUMERIC} & No & \texttt{CHECK (precio >= 0)} & Precio de venta en moneda local (COP). \\
\texttt{activo} & \texttt{BOOLEAN} & No & \texttt{DEFAULT true} & Habilitado en la carta para nuevos pedidos. \\
\midrule
\multicolumn{5}{l}{\textbf{Tabla: public.especialidades} (Autorización de platos para cada cocinero)} \\
\texttt{id} & \texttt{SERIAL} & No & \texttt{PRIMARY KEY} & Identificador de acreditación. \\
\texttt{cocinero\_id} & \texttt{INTEGER} & No & \texttt{FK -> usuarios(id)} & Usuario con rol Cocinero autorizado. \\
\texttt{plato\_id} & \texttt{INTEGER} & No & \texttt{FK -> platos(id)} & Plato que el cocinero sabe elaborar. \\
\bottomrule
\end{tabularx}
\caption{Diccionario de datos: Menú y Especialidades culinarias.}
\label{tab:dict_menu}
\end{table}

\newpage
\subsection{Módulo de Mesas, Reservaciones y Horarios}

\begin{table}[h!]
\centering
\scriptsize
\begin{tabularx}{\textwidth}{l l l p{2.8cm} X}
\toprule
\textbf{Columna} & \textbf{Tipo Dato} & \textbf{Nulo} & \textbf{Claves / Constraints} & \textbf{Descripción Funcional} \\
\midrule
\multicolumn{5}{l}{\textbf{Tabla: public.mesas} (Mesas físicas del establecimiento)} \\
\texttt{id} & \texttt{SERIAL} & No & \texttt{PRIMARY KEY} & Identificador o número físico de la mesa. \\
\texttt{sillas} & \texttt{INTEGER} & No & \texttt{CHECK (sillas > 0)} & Número de puestos o capacidad de comensales. \\
\midrule
\multicolumn{5}{l}{\textbf{Tabla: public.reservaciones} (Tabla particionada de reservaciones solicitadas)} \\
\texttt{id} & \texttt{INTEGER} & No & \texttt{PRIMARY KEY} & Identificador de la reservación. \\
\texttt{cliente\_id} & \texttt{INTEGER} & No & \texttt{FK -> usuarios(id)} & Cliente titular que realiza la reservación. \\
\texttt{cantidad} & \texttt{INTEGER} & No & \texttt{CHECK (cantidad > 0)} & Número de personas que asistirán. \\
\texttt{estado} & \texttt{INTEGER} & No & \texttt{CHECK (1,2,3,4)} & 1=Pendiente, 2=Confirmada, 3=Completada, 4=Cancelada. \\
\midrule
\multicolumn{5}{l}{\textbf{Tabla: public.horarios} (Asignación temporal de mesa a cada reservación)} \\
\texttt{id} & \texttt{SERIAL} & No & \texttt{PRIMARY KEY} & Identificador de agenda. \\
\texttt{mesa\_id} & \texttt{INTEGER} & No & \texttt{FK -> mesas(id)} & Mesa asignada a la reserva. \\
\texttt{reservacion\_id} & \texttt{INTEGER} & No & \texttt{FK -> reservaciones} & Reservación vinculada. \\
\texttt{inicio} & \texttt{TIMESTAMP} & No & -- & Fecha y hora programada de inicio del turno. \\
\texttt{duracion} & \texttt{INTERVAL} & No & \texttt{CHECK (<= '2 hours')} & Duración estimada (fijada en 2 horas). \\
\bottomrule
\end{tabularx}
\caption{Diccionario de datos: Mesas, Reservaciones y Horarios.}
\label{tab:dict_mesas_reservas}
\end{table}

\subsection{Módulo de Pedidos, Órdenes y Preparaciones}

\begin{table}[h!]
\centering
\scriptsize
\begin{tabularx}{\textwidth}{l l l p{2.8cm} X}
\toprule
\textbf{Columna} & \textbf{Tipo Dato} & \textbf{Nulo} & \textbf{Claves / Constraints} & \textbf{Descripción Funcional} \\
\midrule
\multicolumn{5}{l}{\textbf{Tabla: public.pedidos} (Cabecera de comanda en mesa)} \\
\texttt{id} & \texttt{SERIAL} & No & \texttt{PRIMARY KEY} & Identificador del pedido. \\
\texttt{cliente\_id} & \texttt{INTEGER} & No & \texttt{FK -> usuarios(id)} & Cliente que realiza el consumo. \\
\texttt{mesero\_id} & \texttt{INTEGER} & No & \texttt{FK -> usuarios(id)} & Empleado con rol Mesero que atiende. \\
\texttt{mesa\_id} & \texttt{INTEGER} & Sí & \texttt{FK -> mesas(id)} & Mesa donde se sirve el pedido. \\
\texttt{reservacion\_id} & \texttt{INTEGER} & Sí & \texttt{FK -> reservaciones} & Reserva asociada si el pedido deriva de ella. \\
\texttt{creado} & \texttt{TIMESTAMPTZ} & No & \texttt{DEFAULT now()} & Fecha y hora exacta de apertura del pedido. \\
\texttt{token} & \texttt{TEXT} & Sí & \texttt{UNIQUE} & Token de idempotencia contra reenvíos HTTP. \\
\midrule
\multicolumn{5}{l}{\textbf{Tabla: public.ordenes} (Detalle de platos solicitados en un pedido)} \\
\texttt{id} & \texttt{SERIAL} & No & \texttt{PRIMARY KEY} & Identificador de línea de pedido. \\
\texttt{plato\_id} & \texttt{INTEGER} & No & \texttt{FK -> platos(id)} & Ítem culinario ordenado. \\
\texttt{pedido\_id} & \texttt{INTEGER} & No & \texttt{FK -> pedidos(id)} & Pedido cabecera al que pertenece. \\
\texttt{estado} & \texttt{INTEGER} & No & \texttt{CHECK (1,2,3,4)} & 1=Solicitada, 2=En preparación, 4=Preparada, 3=Entregada. \\
\texttt{cantidad} & \texttt{INTEGER} & No & \texttt{CHECK (cantidad > 0)} & Número de porciones solicitadas. \\
\texttt{solicitado} & \texttt{TIMESTAMPTZ} & No & \texttt{DEFAULT now()} & Marca temporal del envío a cocina. \\
\texttt{precio\_unitario} & \texttt{NUMERIC} & Sí & \texttt{CHECK (>= 0)} & Precio congelado al momento del registro. \\
\texttt{preparado} & \texttt{TIMESTAMPTZ} & Sí & -- & Instante en que cocina finaliza el plato. \\
\texttt{entregado} & \texttt{TIMESTAMPTZ} & Sí & -- & Instante en que el mesero sirve en mesa. \\
\midrule
\multicolumn{5}{l}{\textbf{Tabla: public.preparaciones} (Asignación de cocinero a una orden)} \\
\texttt{id} & \texttt{SERIAL} & No & \texttt{PRIMARY KEY} & Identificador de preparación. \\
\texttt{cocinero\_id} & \texttt{INTEGER} & No & \texttt{FK -> usuarios(id)} & Cocinero acreditado que elabora el plato. \\
\texttt{orden\_id} & \texttt{INTEGER} & No & \texttt{FK -> ordenes(id)} & Orden específica asignada. \\
\bottomrule
\end{tabularx}
\caption{Diccionario de datos: Pedidos, Órdenes y Preparaciones.}
\label{tab:dict_pedidos}
\end{table}

\newpage
\section{Arquitectura Física: Tablespaces y Particionamiento}

\subsection{Configuración de Tablespaces en PostgreSQL}
Conforme al Requisito Técnico 6.1, la base de datos aprovecha múltiples espacios de tablas (\textit{tablespaces}) ubicados en rutas independientes del disco duro para segregar físicamente las cargas de entrada/salida (I/O). En el clúster se encuentran configurados los tablespaces de reservaciones:
\begin{itemize}[leftmargin=*]
    \item \texttt{ts\_res\_1}: Directorio físico de almacenamiento para el primer tramo de reservaciones.
    \item \texttt{ts\_res\_2}: Directorio físico de almacenamiento para el segundo tramo de reservaciones.
    \item \texttt{ts\_res\_3}: Directorio físico de almacenamiento para el tercer tramo de reservaciones históricas.
\end{itemize}

Esta arquitectura permite separar los discos mecánicos o volúmenes SSD, balanceando la concurrencia entre transacciones activas de lectura y escrituras analíticas.

\subsection{Estrategia de Particionamiento de Tablas}
La tabla \texttt{reservaciones} se encuentra particionada por rangos (\texttt{PARTITION BY RANGE (id)}), distribuyendo sus registros en particiones hijas alojadas en los tablespaces independientes:
\begin{lstlisting}[style=sqlstyle]
-- Tabla principal declarativa particionada
CREATE TABLE public.reservaciones (
    id integer NOT NULL,
    cliente_id integer NOT NULL,
    cantidad integer NOT NULL,
    estado integer NOT NULL,
    CONSTRAINT ck_reservaciones_cantidad_positiva CHECK (cantidad > 0),
    CONSTRAINT ck_reservaciones_estado CHECK (estado IN (1, 2, 3, 4))
) PARTITION BY RANGE (id);

-- Asignación de particiones a sus respectivos tablespaces
CREATE TABLE public.reservaciones_1 PARTITION OF public.reservaciones
    FOR VALUES FROM (1) TO (15001) TABLESPACE ts_res_1;

CREATE TABLE public.reservaciones_2 PARTITION OF public.reservaciones
    FOR VALUES FROM (15001) TO (30001) TABLESPACE ts_res_2;

CREATE TABLE public.reservaciones_3 PARTITION OF public.reservaciones
    FOR VALUES FROM (30001) TO (45001) TABLESPACE ts_res_3;

-- Partición por defecto para cubrir crecimientos futuros sin bloqueo
CREATE TABLE public.reservaciones_futuras PARTITION OF public.reservaciones DEFAULT;
\end{lstlisting}

\paragraph{Particionamiento por Fechas (Diseño de Eventos):}
Para tablas de series temporales de alta cardinalidad como \texttt{ordenes} y \texttt{pedidos}, el particionamiento idóneo sobre los tablespaces definidos se formula mediante rangos de fechas (\texttt{RANGE (solicitado)}), permitiendo la exclusión de particiones (\textit{Partition Pruning}) en consultas mensuales o anuales:
\begin{lstlisting}[style=sqlstyle]
-- Ejemplo de particionamiento temporal por meses para ventas
CREATE TABLE public.ordenes_part (
    id SERIAL,
    solicitado TIMESTAMPTZ NOT NULL,
    ...
) PARTITION BY RANGE (solicitado);

CREATE TABLE public.ordenes_2026_q3 PARTITION OF public.ordenes_part
    FOR VALUES FROM ('2026-07-01') TO ('2026-10-01') TABLESPACE ts_res_2;

CREATE TABLE public.ordenes_2026_q4 PARTITION OF public.ordenes_part
    FOR VALUES FROM ('2026-10-01') TO ('2027-01-01') TABLESPACE ts_res_1;
\end{lstlisting}

\newpage
\section{Integridad de Datos y Disparadores (Triggers)}

El Requisito Técnico 6.3 exige la implementación de mecanismos de integridad relacional activa en el motor de base de datos. En el script \texttt{01\_integridad\_y\_triggers.sql} se implementaron cinco disparadores (\textit{triggers}) clave en PL/pgSQL:

\subsection{Trigger 1: Prevención de Cruces de Reservaciones (2 Horas)}
Evita matemáticamente que se solapen dos reservaciones no canceladas en la misma mesa durante un bloque mínimo de dos horas:
\begin{lstlisting}[style=sqlstyle]
CREATE OR REPLACE FUNCTION public.fn_trg_validar_cruce_reserva()
RETURNS TRIGGER AS $$
DECLARE
    v_capacidad_mesa integer;
    v_personas_reserva integer;
    v_estado_reserva integer;
    v_conflicto_id integer;
    v_duracion interval;
BEGIN
    SELECT cantidad, estado INTO v_personas_reserva, v_estado_reserva
    FROM public.reservaciones WHERE id = NEW.reservacion_id;

    IF v_estado_reserva = 4 THEN RETURN NEW; END IF; -- Cancelada

    SELECT sillas INTO v_capacidad_mesa FROM public.mesas WHERE id = NEW.mesa_id;
    IF v_personas_reserva > v_capacidad_mesa THEN
        RAISE EXCEPTION 'Capacidad excedida: Reserva para % personas en mesa de % sillas.',
            v_personas_reserva, v_capacidad_mesa;
    END IF;

    v_duracion := GREATEST(COALESCE(NEW.duracion, interval '2 hours'), interval '2 hours');

    SELECT h.id INTO v_conflicto_id
    FROM public.horarios h
    JOIN public.reservaciones r ON r.id = h.reservacion_id
    WHERE h.mesa_id = NEW.mesa_id AND r.estado <> 4
      AND h.id <> COALESCE(NEW.id, -1) AND h.reservacion_id <> NEW.reservacion_id
      AND tsrange(h.inicio, h.inicio + GREATEST(h.duracion, interval '2 hours'), '[)')
       && tsrange(NEW.inicio, NEW.inicio + v_duracion, '[)')
    LIMIT 1;

    IF v_conflicto_id IS NOT NULL THEN
        RAISE EXCEPTION 'Cruce de horario: Mesa #% ya reservada (conflicto #%).',
            NEW.mesa_id, v_conflicto_id;
    END IF;
    RETURN NEW;
END;
$$ LANGUAGE plpgsql;

CREATE TRIGGER trg_validar_cruce_reserva
BEFORE INSERT OR UPDATE OF mesa_id, inicio, duracion ON public.horarios
FOR EACH ROW EXECUTE FUNCTION public.fn_trg_validar_cruce_reserva();
\end{lstlisting}

\subsection{Trigger 2: Validación Estricta de Cupo Total del Restaurante}
Garantiza que la sumatoria de comensales concurrentes en la ventana horaria no supere el cupo total del restaurante ($\sum \text{sillas}$):
\begin{lstlisting}[style=sqlstyle]
CREATE OR REPLACE FUNCTION public.fn_trg_validar_cupo_restaurante()
RETURNS TRIGGER AS $$
DECLARE
    v_cupo_total integer;
    v_ocupacion_pico integer;
    v_personas_nuevas integer;
BEGIN
    SELECT cantidad INTO v_personas_nuevas FROM public.reservaciones WHERE id = NEW.reservacion_id;
    SELECT COALESCE(SUM(sillas), 0) INTO v_cupo_total FROM public.mesas;

    WITH ocupaciones AS (
        SELECT h.inicio, h.inicio + GREATEST(h.duracion, interval '2 hours') AS fin, r.cantidad
        FROM public.horarios h JOIN public.reservaciones r ON r.id = h.reservacion_id
        WHERE r.estado <> 4 AND h.reservacion_id <> NEW.reservacion_id
          AND h.inicio < NEW.inicio + interval '2 hours'
          AND h.inicio + GREATEST(h.duracion, interval '2 hours') > NEW.inicio
    ),
    puntos AS (
        SELECT NEW.inicio AS t UNION SELECT inicio FROM ocupaciones WHERE inicio >= NEW.inicio
    )
    SELECT COALESCE(MAX((SELECT COALESCE(SUM(cantidad),0) FROM ocupaciones o WHERE o.inicio<=p.t AND o.fin>p.t)),0)
    INTO v_ocupacion_pico FROM puntos p;

    IF (v_ocupacion_pico + v_personas_nuevas) > v_cupo_total THEN
        RAISE EXCEPTION 'Cupo total superado: Capacidad=%, concurrencia solicitada=%.',
            v_cupo_total, (v_ocupacion_pico + v_personas_nuevas);
    END IF;
    RETURN NEW;
END;
$$ LANGUAGE plpgsql;

CREATE TRIGGER trg_validar_cupo_restaurante
BEFORE INSERT OR UPDATE OF inicio, duracion ON public.horarios
FOR EACH ROW EXECUTE FUNCTION public.fn_trg_validar_cupo_restaurante();
\end{lstlisting}

\subsection{Trigger 3: Acreditación de Especialidad al Cocinero}
Comprueba que el cocinero que toma una orden posea el rol 'Cocinero' y tenga registrada la especialidad del plato correspondiente:
\begin{lstlisting}[style=sqlstyle]
CREATE OR REPLACE FUNCTION public.fn_trg_validar_preparacion_cocinero()
RETURNS TRIGGER AS $$
DECLARE
    v_es_cocinero boolean;
    v_tiene_especialidad boolean;
    v_plato_id integer;
BEGIN
    SELECT EXISTS (
        SELECT 1 FROM public.actuaciones a JOIN public.roles r ON r.id = a.rol_id
        WHERE a.usuario_id = NEW.cocinero_id AND r.nombre = 'Cocinero'
    ) INTO v_es_cocinero;

    IF NOT v_es_cocinero THEN
        RAISE EXCEPTION 'El usuario #% no tiene el rol de Cocinero.', NEW.cocinero_id;
    END IF;

    SELECT plato_id INTO v_plato_id FROM public.ordenes WHERE id = NEW.orden_id;

    SELECT EXISTS (
        SELECT 1 FROM public.especialidades
        WHERE cocinero_id = NEW.cocinero_id AND plato_id = v_plato_id
    ) INTO v_tiene_especialidad;

    IF NOT v_tiene_especialidad THEN
        RAISE EXCEPTION 'Especialidad no autorizada: Cocinero #% no autorizado para plato #%.',
            NEW.cocinero_id, v_plato_id;
    END IF;
    RETURN NEW;
END;
$$ LANGUAGE plpgsql;

CREATE TRIGGER trg_validar_preparacion_cocinero
BEFORE INSERT OR UPDATE OF cocinero_id, orden_id ON public.preparaciones
FOR EACH ROW EXECUTE FUNCTION public.fn_trg_validar_preparacion_cocinero();
\end{lstlisting}

\newpage
\section{Consultas SQL de los Casos de Uso y Reportes}

El script \texttt{02\_consultas\_casos\_de\_uso.sql} consolida las consultas relacionales requeridas. A continuación se presentan las más representativas:

\subsection{CU-01: Cupo Total del Restaurante y Aforo}
\begin{lstlisting}[style=sqlstyle]
SELECT
    count(id)               AS total_mesas,
    sum(sillas)             AS cupo_total_restaurante,
    min(sillas)             AS capacidad_minima_mesa,
    max(sillas)             AS capacidad_maxima_mesa,
    round(avg(sillas), 2)   AS promedio_sillas_por_mesa
FROM public.mesas;
\end{lstlisting}

\subsection{CU-02: Disponibilidad y Alerta de Reservas para el Maitre}
\begin{lstlisting}[style=sqlstyle]
-- Alertas de reservas proximas a comenzar en ventana <= 2 horas
SELECT
    r.id AS reservacion_id,
    u.nombre AS cliente,
    h.mesa_id,
    h.inicio,
    (h.inicio - LOCALTIMESTAMP) AS tiempo_restante,
    r.cantidad AS personas
FROM public.reservaciones r
JOIN public.usuarios u ON u.id = r.cliente_id
JOIN public.horarios h ON h.reservacion_id = r.id
WHERE r.estado IN (1, 2)
  AND h.inicio >= LOCALTIMESTAMP
  AND h.inicio <= LOCALTIMESTAMP + interval '2 hours'
ORDER BY h.inicio ASC;
\end{lstlisting}

\subsection{CU-05: Cola de Cocina Priorizada por Tiempo Descendente}
El cocinero consulta únicamente las órdenes pendientes que sabe cocinar, priorizando los platos más demorados de primero:
\begin{lstlisting}[style=sqlstyle]
SELECT
    o.id AS orden_id,
    ped.mesa_id,
    p.nombre AS plato,
    p.tiempo AS tiempo_preparacion,
    o.cantidad,
    o.solicitado
FROM public.ordenes o
JOIN public.platos p ON p.id = o.plato_id
JOIN public.pedidos ped ON ped.id = o.pedido_id
JOIN public.especialidades esp ON esp.plato_id = p.id AND esp.cocinero_id = :p_cocinero_id
WHERE o.estado = 1 -- Solicitadas y disponibles
ORDER BY p.tiempo DESC, o.solicitado ASC;
\end{lstlisting}

\subsection{CU-07: Reportes Estadísticos Analíticos}
\begin{lstlisting}[style=sqlstyle]
-- Reporte de Ventas por periodo mensual (date_trunc)
SELECT
    date_trunc('month', COALESCE(o.entregado, o.solicitado) AT TIME ZONE 'America/Bogota')::date AS periodo,
    count(DISTINCT o.pedido_id) AS total_pedidos,
    sum(o.cantidad) AS platos_servidos,
    sum(o.cantidad * o.precio_unitario) FILTER (WHERE o.entregado IS NOT NULL) AS ventas_verificadas,
    sum(o.cantidad * COALESCE(o.precio_unitario, p.precio)) AS facturacion_total
FROM public.ordenes o
JOIN public.platos p ON p.id = o.plato_id
WHERE o.estado = 3 -- Entregado
GROUP BY 1 ORDER BY 1;

-- Ranking Top de Platos mas solicitados
SELECT
    p.nombre AS plato,
    t.nombre AS categoria,
    sum(o.cantidad) AS unidades_vendidas,
    sum(o.cantidad * COALESCE(o.precio_unitario, p.precio)) AS facturacion
FROM public.ordenes o
JOIN public.platos p ON p.id = o.plato_id
JOIN public.tipos t ON t.id = p.tipo_id
GROUP BY p.id, p.nombre, t.nombre
ORDER BY unidades_vendidas DESC LIMIT 10;
\end{lstlisting}

\newpage
\section{Estrategia de Pruebas y Validación}

\subsection{Conjunto de Datos de Prueba Determinista}
Para garantizar la reproducibilidad científica de las evaluaciones, se utilizó la semilla determinista \texttt{03\_datos\_prueba\_Garcia\_Pico.sql} con factor \texttt{data\_seed = 42}. Esta semilla genera:
\begin{itemize}
    \item 32 mesas con configuraciones variadas de 2 a 8 sillas.
    \item Plantilla completa de empleados (4 Maitres, 12 Cocineros con capacitación cruzada, 10 Meseros y 2 Administradores).
    \item Catálogo completo de 24 platos clasificados en Entradas, Fuertes, Postres y Bebidas.
    \item 10.000 clientes sintéticos con más de 18.000 reservaciones y órdenes históricas.
\end{itemize}

\subsection{Pruebas Automatizadas de Integración}
La suite de pruebas disponible en \texttt{tests/integration.php} valida exhaustivamente:
\begin{enumerate}
    \item \textbf{Aislamiento de Entorno:} Restricción estricta sobre la base de pruebas \texttt{restaurante\_test} para impedir mutación accidental de datos en producción.
    \item \textbf{Idempotencia:} Uso de advisory locks (\texttt{pg\_advisory\_xact\_lock}) y tokens de un solo uso que rechazan duplicación de pedidos ante reenvíos concurrentes.
    \item \textbf{Control de Estados y Privilegios:} Imposibilidad de que un mesero entregue un pedido antes de ser preparado, o que un cocinero sin especialidad tome una comanda.
\end{enumerate}

\newpage
\section{Guía y Libreto de Sustentación Oral (Rúbrica 20\%)}

Para la demostración en vivo ante el docente titular, se recomienda seguir el siguiente guion estructurado:

\begin{enumerate}[leftmargin=*]
    \item \textbf{Introducción y Modelo Relacional (3 min):}
    \begin{itemize}
        \item Proyectar el diagrama ER destacando la entidad asociativa \texttt{actuaciones} que desacopla los usuarios de sus roles, y la entidad asociativa \texttt{especialidades} que enlaza a cocineros con platos autorizados.
        \item Demostrar el diccionario de datos y las restricciones de no nulidad y valores positivos.
    \end{itemize}
    \item \textbf{Tablespaces y Particionamiento Físico (2 min):}
    \begin{itemize}
        \item Ejecutar en terminal: \texttt{\textbackslash db+} en \texttt{psql} para evidenciar los tablespaces creados en directorios separados del disco duro.
        \item Consultar el catálogo: \texttt{SELECT relname, relpartbound FROM pg\_class WHERE relispartition = true;} para mostrar las particiones activas de \texttt{reservaciones} y la partición predeterminada.
    \end{itemize}
    \item \textbf{Demostración de Disparadores de Integridad (3 min):}
    \begin{itemize}
        \item \textit{Caso Cruce de Horario:} Intentar insertar un horario sobre una mesa que ya esté reservada en la misma ventana de 2 horas. Mostrar el mensaje de error provocado por el trigger \texttt{trg\_validar\_cruce\_reserva}.
        \item \textit{Caso Especialidad Culinaria:} Intentar asignar en \texttt{preparaciones} a un cocinero sobre un plato que no tiene registrado en \texttt{especialidades}. Mostrar la excepción del trigger \texttt{trg\_validar\_preparacion\_cocinero}.
    \end{itemize}
    \item \textbf{Ejecución de Casos de Uso y Reportes (4 min):}
    \begin{itemize}
        \item Ejecutar la consulta del CU-05 mostrando cómo la cola de cocina se ordena estrictamente por tiempo descendente de preparación.
        \item Ejecutar el reporte de ventas con \texttt{date\_trunc('month', ...)} y el ranking de platos más vendidos.
    \end{itemize}
\end{enumerate}

\section{Declaración de Uso de Inteligencia Artificial}
En estricto cumplimiento del \textbf{Requisito 9} de la especificación del proyecto, se declara que para el análisis de requisitos, diseño de scripts modulares, estructuración del informe formal en \LaTeX\ y codificación colaborativa se utilizaron herramientas de inteligencia artificial de \textit{live coding} (Antigravity AI / Google DeepMind).

El estudiante asume la total autoría y responsabilidad intelectual y técnica sobre el código, los disparadores, las consultas SQL y el documento entregado, habiendo validado y verificado cada una de las pruebas y encontrándose plenamente capacitado para su sustentación técnica oral ante el docente.

\end{document}

